\begin{proof}[Proof of \cref{obj:lem:frontier-testing-priors}]
\begin{enumerate}
\item Define the three candidate scales and the separation constant by
\[
a_n:=n^{-1/4},\qquad
b_n:=(n^2\epsilon)^{-1/7},\qquad
c_n:=(n\epsilon)^{-1/2},\qquad
\kappa:=2^{-12}.
\]
All three scales are positive. Set
\[
h_{\mathrm L}:=\frac{r(n,\epsilon)}4,\qquad
t:=\kappa h_{\mathrm L},\qquad
\delta_{\mathrm L}:=h_{\mathrm L}^{5},\qquad
\Delta:=t.
\]
Since \(0<r(n,\epsilon)\le1\),
\[
0<h_{\mathrm L}\le\frac14,
\qquad
\Delta=\frac{\kappa}{4}r(n,\epsilon)
      =2^{-14}r(n,\epsilon).
\]
We now split according to the sign-family condition in the statement.
% lean: CausalSmith.Stat.PrivateCateRoughdesign.frontier_testing_priors

\item Suppose that
\[
1<n\epsilon
\quad\text{and}\quad
\bigl[(b_n\le a_n\ \text{and}\ c_n\le a_n)
\ \text{or}\
(a_n<b_n\ \text{and}\ c_n<b_n)\bigr].
\]
At radius \(h_{\mathrm L}\), let the active index set and its binary sign vectors be
\[
\begin{aligned}
\mathcal J_{h_{\mathrm L}}
:=\biggl\{j\in\mathbb Z:\;&
\left\lfloor\frac{x_0-h_{\mathrm L}}{\delta_{\mathrm L}}\right\rfloor-1
\le j\le
\left\lceil\frac{x_0+h_{\mathrm L}}{\delta_{\mathrm L}}\right\rceil+1,\\
&j\delta_{\mathrm L}-\delta_{\mathrm L}<x_0+h_{\mathrm L},
\quad
x_0-h_{\mathrm L}<j\delta_{\mathrm L}+\delta_{\mathrm L}
\biggr\},
\qquad x_0=\frac12,\\
\Lambda_{h_{\mathrm L}}&:=\{\lambda:\mathcal J_{h_{\mathrm L}}\to\{0,1\}\}.
\end{aligned}
\]
The integer interval in this definition is bounded, so \(\mathcal J_{h_{\mathrm L}}\) is finite. Consequently,
\[
|\Lambda_{h_{\mathrm L}}|=2^{|\mathcal J_{h_{\mathrm L}}|}>0.
\]
This also covers an empty active set, for which there is one empty sign vector.

Choose the singleton null family \((P_0)\) and the alternative family
\((P_\lambda)_{\lambda\in\Lambda_{h_{\mathrm L}}}\) from
\cref{obj:def:cosine-family}. By
\cref{obj:lem:positive-family-certificate},
\[
P_0\in\mathcal P,\qquad \theta(P_0)=0,\qquad
P_\lambda\in\mathcal P,\qquad \theta(P_\lambda)=t=\Delta
\quad(\lambda\in\Lambda_{h_{\mathrm L}}).
\]
Their uniform product-law averages are
\[
\begin{aligned}
Q^{(0)}&=(P_0^{\mathrm{obs}})^{\otimes n},\\
Q^{(1)}
&=\frac1{|\Lambda_{h_{\mathrm L}}|}
  \sum_{\lambda\in\Lambda_{h_{\mathrm L}}}
  (P_\lambda^{\mathrm{obs}})^{\otimes n}\\
&=2^{-|\mathcal J_{h_{\mathrm L}}|}
  \sum_{\lambda\in\Lambda_{h_{\mathrm L}}}
  (P_\lambda^{\mathrm{obs}})^{\otimes n}
 =Q_n.
\end{aligned}
\]
The equality follows by enumerating every binary sign vector exactly once.
% lean: CausalSmith.Stat.PrivateCateRoughdesign.finiteProductMixture_signs

\item First consider the sampling subcase
\[
b_n\le a_n,\qquad c_n\le a_n.
\]
Since \(n\ge2\), we have \(a_n\le1\), and hence
\[
r(n,\epsilon)=a_n,\qquad
h_{\mathrm L}=\frac{a_n}{4}.
\]
The power identity \(a_n^8=n^{-2}\) gives
\[
n^2h_{\mathrm L}^{8}=\frac1{4^8}.
\]
Define the nonnegative exponent
\[
\eta_{\mathrm{samp}}
:=160n^2t^2h_{\mathrm L}\delta_{\mathrm L}
 =160\kappa^2n^2h_{\mathrm L}^{8}.
\]
It satisfies
\[
0\le\eta_{\mathrm{samp}}
 =160\cdot2^{-40}\le\frac1{16}.
\]
The chi-square estimate of
\cref{obj:lem:positive-family-certificate} uses the bounded-differences
exponential-moment conclusion of
\citep[Theorem~6.2, displayed log-mgf conclusion in its proof, p.~167]{BoucheronLugosiMassart2013}.
It supplies
\[
\chi^2(Q_n,Q^{(0)})\le e^{\eta_{\mathrm{samp}}}-1,
\]
together with absolute continuity and square-integrability of the centered
likelihood ratio.

Convexity of the exponential between \(0\) and \(1\), and \(e<3\), imply
\[
e^{\eta_{\mathrm{samp}}}-1
\le\eta_{\mathrm{samp}}(e-1)
\le2\eta_{\mathrm{samp}}.
\]
Therefore,
\[
\chi^2(Q_n,Q^{(0)})\le2\eta_{\mathrm{samp}}\le\frac14,
\qquad
\sqrt{\chi^2(Q_n,Q^{(0)})}\le\frac12.
\]
For every allowed \(M\), symmetry and the two comparison inequalities in
\cref{obj:lem:private-kernel-comparison} now give
\[
d_{\mathrm{TV}}(MQ^{(0)},MQ^{(1)})
=d_{\mathrm{TV}}(MQ_n,MQ^{(0)})
\le d_{\mathrm{TV}}(Q_n,Q^{(0)})
\le\frac12\sqrt{\chi^2(Q_n,Q^{(0)})}
\le\frac14.
\]
% lean: CausalSmith.Stat.PrivateCateRoughdesign.sampling_sign_output_TV_le

\item Next consider the sparse subcase
\[
a_n<b_n,\qquad c_n<b_n.
\]
Here \(n\epsilon>1\), so \(n^2\epsilon>1\) and \(0<b_n\le1\). Thus
\[
r(n,\epsilon)=b_n,\qquad h_{\mathrm L}=\frac{b_n}{4}.
\]
To establish the occupancy condition, apply the increasing logarithm to
\(c_n\le b_n\). This gives
\[
-\frac12(\log n+\log\epsilon)
\le-\frac17(2\log n+\log\epsilon),
\qquad
3\log n+5\log\epsilon\ge0.
\]
Consequently,
\[
\log(nb_n^5)
=\log n-\frac57(2\log n+\log\epsilon)
=-\frac17(3\log n+5\log\epsilon)
\le0.
\]
It follows that
\[
nb_n^5\le1,\qquad
n\delta_{\mathrm L}
=\frac{nb_n^5}{4^5}
\le\frac1{1024}\le\frac1{128}.
\]

The component estimate in
\cref{obj:lem:measurable-component-contraction} uses Cayley's labeled-tree
enumeration from
\citep[Section~5.6, Theorem~5.40]{KellerTrotterAppliedCombinatorics}.
With the occupancy condition just established, that estimate supplies
\[
d_{\mathrm{TV}}(MQ_n,MQ^{(0)})
\le2048\epsilon t n^2h_{\mathrm L}\delta_{\mathrm L}.
\]
The identity
\[
b_n^7=(n^2\epsilon)^{-1}
\]
calibrates its right-hand side:
\[
\begin{aligned}
2048\epsilon t n^2h_{\mathrm L}\delta_{\mathrm L}
&=2048\epsilon\kappa n^2h_{\mathrm L}^{7}\\
&=\frac{2048\kappa}{4^7}(n^2\epsilon)b_n^7\\
&=\frac{2048\kappa}{4^7}
 =2^{-15}\le\frac14.
\end{aligned}
\]
Symmetry therefore yields the required bound for
\(d_{\mathrm{TV}}(MQ^{(0)},MQ^{(1)})\).
% lean: CausalSmith.Stat.PrivateCateRoughdesign.measurable_component_contraction

\item Suppose now that the sign-family condition fails. Choose the singleton
families \((P_0)\) and \((P_1)\), with \(P_1\) defined at radius
\(h_{\mathrm L}\) by \cref{obj:def:direct-family}. The null membership and
target follow from \cref{obj:lem:positive-family-certificate}. We verify
the corresponding properties of \(P_1\).

For the localized profile \(q\) in \cref{obj:def:direct-family},
\[
0\le q(x)\le1,\qquad q(x_0)=1,\qquad x_0=\frac12.
\]
The calibration gives
\[
0<t=\kappa h_{\mathrm L}\le2^{-14}<\frac12,
\qquad
\frac12\le\mu_{1,P_1}(x)=\frac12+tq(x)\le1.
\]
To check the contrast modulus, define, for \(x\in[0,1]\),
\[
g_{\mathrm{loc}}(x)
:=\cos\!\left(
\frac{\pi}{2h_{\mathrm L}}
\min\{h_{\mathrm L},\max\{-h_{\mathrm L},x-x_0\}\}
\right).
\]
At and beyond either support boundary the cosine is zero, so
\[
q(x)=g_{\mathrm{loc}}(x)^4,\qquad |g_{\mathrm{loc}}(x)|\le1.
\]
Clipping to an interval and the cosine function are both \(1\)-Lipschitz;
therefore
\[
|g_{\mathrm{loc}}(x)-g_{\mathrm{loc}}(x')|
\le\frac{\pi}{2h_{\mathrm L}}|x-x'|.
\]
Factoring a difference of squares twice gives
\[
\begin{aligned}
|g_{\mathrm{loc}}(x)^2-g_{\mathrm{loc}}(x')^2|
&\le2|g_{\mathrm{loc}}(x)-g_{\mathrm{loc}}(x')|,\\
|q(x)-q(x')|
&\le2|g_{\mathrm{loc}}(x)^2-g_{\mathrm{loc}}(x')^2|
 \le\frac{2\pi}{h_{\mathrm L}}|x-x'|.
\end{aligned}
\]
Hence
\[
|\tau_{P_1}(x)-\tau_{P_1}(x')|
=t|q(x)-q(x')|
\le2\pi\kappa|x-x'|
\le L|x-x'|,
\]
because \(\pi<4\), \(\kappa=2^{-12}\), and \(L=3\) by
\cref{obj:synth_4}.

The covariate density is one, and the propensity and control regression
are constant \(1/2\). In particular,
\[
|e_{P_1}(x)-e_{P_1}(x')|
=|\mu_{0,P_1}(x)-\mu_{0,P_1}(x')|
=0\le L|x-x'|^{1/10}.
\]
These facts give the density, overlap, and nuisance smoothness conditions.
The conditional independent generation in
\cref{obj:def:direct-family} gives exchangeability, and setting
\(Y=Y(A)\) gives consistency. Thus, using
\cref{obj:def:complete-model},
\[
P_1\in\mathcal P,\qquad
\theta(P_1)=tq(x_0)=t=\Delta.
\]
Uniform averaging over either singleton leaves its product law unchanged:
\[
Q^{(0)}=(P_0^{\mathrm{obs}})^{\otimes n},\qquad
Q^{(1)}=(P_1^{\mathrm{obs}})^{\otimes n}.
\]
% lean: CausalSmith.Stat.PrivateCateRoughdesign.directAlternative_completeModel

\item Construct the direct dataset coupling as follows. For each
\(x\in[0,1]\), define the crossing mass
\[
\omega_{\mathrm{cross}}(x):=\frac{tq(x)}2,
\qquad
0\le\omega_{\mathrm{cross}}(x)\le\frac14.
\]
The conditional coupling has the following five atoms; every other pair
has mass zero:
\[
\begin{array}{c|c|c}
\text{First observed record}&\text{Second observed record}&\text{Mass}\\ \hline
(x,0,0)&(x,0,0)&1/4\\
(x,0,1)&(x,0,1)&1/4\\
(x,1,0)&(x,1,0)&1/4-\omega_{\mathrm{cross}}(x)\\
(x,1,1)&(x,1,1)&1/4\\
(x,1,0)&(x,1,1)&\omega_{\mathrm{cross}}(x)
\end{array}
\]
The masses are nonnegative and sum to one:
\[
\frac14+\frac14+
\left(\frac14-\omega_{\mathrm{cross}}(x)\right)
+\frac14+\omega_{\mathrm{cross}}(x)=1.
\]
The first marginal has mass \(1/4\) at each binary treatment--outcome
configuration. In the second marginal the control-cell masses are
\(1/4,1/4\), while the treated-cell masses satisfy
\[
\frac14-\omega_{\mathrm{cross}}(x)
=\frac12\left(\frac12-tq(x)\right),\qquad
\frac14+\omega_{\mathrm{cross}}(x)
=\frac12\left(\frac12+tq(x)\right).
\]
Thus these marginals are exactly the conditional observed laws of
\(P_0\) and \(P_1\).

Integrate this measurable conditional coupling over uniform \(X\), and
take \(n\) independent copies. Group their first and second coordinates
into \(D,D'\in\mathcal O^n\), and define
\[
\gamma_{\mathrm{dir}}:=\mathcal L(D,D').
\]
The product construction gives
\[
\mathcal L(D)=Q^{(0)},\qquad \mathcal L(D')=Q^{(1)}.
\]
At a fixed covariate, precisely the crossing atom contributes to record
disagreement. Since \(q\le1\) and vanishes when
\(|x-x_0|\ge h_{\mathrm L}\),
\[
\int_0^1q(x)\,dx
\le\int_0^1\mathbf1\{|x-x_0|<h_{\mathrm L}\}\,dx
\le2h_{\mathrm L}.
\]
Adding the \(n\) record disagreement costs therefore yields
\[
\begin{aligned}
\int d_{\mathrm H}(D,D')\,\gamma_{\mathrm{dir}}(dD,dD')
&=\sum_{i=1}^n\Pr\{D_i\ne D'_i\}\\
&=n\int_0^1\omega_{\mathrm{cross}}(x)\,dx\\
&=\frac{nt}{2}\int_0^1q(x)\,dx
\le nth_{\mathrm L}.
\end{aligned}
\]
By \cref{obj:lem:private-kernel-comparison} and the definition of the
infimum coupling cost,
\[
\begin{aligned}
d_{\mathrm{TV}}(MQ^{(0)},MQ^{(1)})
&\le2\epsilon W_{\mathrm H}(Q^{(0)},Q^{(1)})\\
&\le2\epsilon
 \int d_{\mathrm H}(D,D')\,\gamma_{\mathrm{dir}}(dD,dD')\\
&\le2\epsilon nth_{\mathrm L}.
\end{aligned}
\]
% lean: CausalSmith.Stat.PrivateCateRoughdesign.directDatasetCoupling_cost_le

\item It remains to calibrate this direct privacy bound. We first show
that failure of the sign-family condition implies
\[
r(n,\epsilon)=\min\{1,c_n\}.
\]
If \(n\epsilon>1\), suppose that \(c_n<a_n\). When \(b_n\le a_n\), the
sampling sign condition holds. When \(a_n<b_n\), we also have
\(c_n<b_n\), so the sparse sign condition holds. Both contradict the
present branch. Thus \(a_n\le c_n\). If \(c_n<b_n\), then
\(a_n\le c_n<b_n\), again giving the sparse sign condition. Consequently,
\[
a_n\le c_n,\qquad b_n\le c_n,\qquad
r(n,\epsilon)=\min\{1,c_n\}.
\]
If \(n\epsilon\le1\), monotonicity of the negative power gives
\[
c_n=(n\epsilon)^{-1/2}\ge1,\qquad
r(n,\epsilon)=1=\min\{1,c_n\}.
\]
This proves the identity throughout the direct branch, including
\(n\epsilon=1\) and ties among the candidate scales.

Since all quantities are nonnegative,
\[
r(n,\epsilon)\le c_n,\qquad
c_n^2=(n\epsilon)^{-1},\qquad
r(n,\epsilon)^2\le(n\epsilon)^{-1},
\qquad
(n\epsilon)r(n,\epsilon)^2\le1.
\]
Using \(t=\kappa h_{\mathrm L}\) and
\(h_{\mathrm L}=r(n,\epsilon)/4\), we obtain
\[
2\epsilon nth_{\mathrm L}
=2\epsilon n\kappa h_{\mathrm L}^2
=\frac{\kappa}{8}(n\epsilon)r(n,\epsilon)^2
\le\frac{\kappa}{8}.
\]
For every standard Borel output space and every private kernel in the
statement, the direct comparison therefore satisfies
\[
d_{\mathrm{TV}}(MQ^{(0)},MQ^{(1)})
\le\frac{\kappa}{8}
=2^{-15}\le\frac14.
\]
Together with the two sign subcases, this completes the construction and
the claimed output comparison.
% lean: CausalSmith.Stat.PrivateCateRoughdesign.testing_direct_privacy_scale_le
\end{enumerate}
\end{proof}
